% ============================================================================
% APPENDIX D - MATHEMATICAL PROOFS
% ============================================================================

\labelsec{appendix_proofs}

This appendix provides formal proofs of key mathematical properties.

\section{D.1 Proof that $\eta_r \in [0,1]$}

\begin{theorem}
The Shambhian index $\eta_r = 1 - \lambda_2/\lambda_1$ satisfies $0 \leq \eta_r \leq 1$ for all non-degenerate distributions.
\end{theorem}

\begin{proof}
By definition, $\lambda_1$ and $\lambda_2$ are eigenvalues of the inertia tensor $I$:

\begin{align}
    I = \begin{pmatrix} I_{xx} & I_{xy} \\ I_{xy} & I_{yy} \end{pmatrix}
\end{align}

Properties of positive semi-definite matrices:
\begin{enumerate}
    \item Both eigenvalues are non-negative: $\lambda_1 \geq 0, \lambda_2 \geq 0$
    \item We order them: $\lambda_1 \geq \lambda_2$
\end{enumerate}

Lower bound: $\eta_r = 1 - \lambda_2/\lambda_1$

Since $\lambda_2 \leq \lambda_1$, we have $\lambda_2/\lambda_1 \leq 1$, thus:
\begin{align}
    \eta_r = 1 - \lambda_2/\lambda_1 \geq 1 - 1 = 0
\end{align}

Upper bound: Since $\lambda_2 \geq 0$, we have $\lambda_2/\lambda_1 \geq 0$, thus:
\begin{align}
    \eta_r = 1 - \lambda_2/\lambda_1 \leq 1 - 0 = 1
\end{align}

Therefore, $\eta_r \in [0, 1]$. $\square$

\end{proof}

\section{D.2 Connection to Classical Skewness}

\begin{theorem}
The Mamton index relates to classical univariate skewness as follows:

\begin{align}
    \Omega = \norm{\vec{S}} \propto \left( \gamma_x^2 + \gamma_y^2 \right)^{1/2}
\end{align}

where $\gamma_x, \gamma_y$ are classical skewness coefficients.
\end{theorem}

\begin{proof}
Classical centered skewness: $\gamma_x = E[(X - \mu_x)^3]/\sigma_x^3$

Our raw skewness components:
\begin{align}
    S_{30} &= \sum_i (x_i - \bar{x})^3 \propto \gamma_x \sigma_x^3 \\
    S_{03} &= \sum_i (y_i - \bar{y})^3 \propto \gamma_y \sigma_y^3
\end{align}

Therefore:
\begin{align}
    \Omega &= \sqrt{S_{30}^2 + S_{03}^2 + S_{21}^2 + S_{12}^2} \\
    &\propto \sqrt{\gamma_x^2 \sigma_x^6 + \gamma_y^2 \sigma_y^6 + \ldots}
\end{align}

For isotropic distributions ($\sigma_x \approx \sigma_y$, $\gamma_x \approx \gamma_y$):
\begin{align}
    \Omega \propto \sqrt{\gamma_x^2 + \gamma_y^2}
\end{align}

$\square$
\end{proof}

\section{D.3 Coordinate-Frame Independence}

\begin{theorem}
The geometric unit values are invariant under rotations of the coordinate frame.
\end{theorem}

\begin{proof}
We show this for $v_0$; similar arguments apply to $\eta_r$ and $\Omega$.

The curvature of a curve is a geometric property independent of coordinate system choice. If $\vec{C}(h) = [x(h), y(h), h]$, then:

\begin{align}
    v_0 &= \norm{\frac{d^2\vec{C}}{dh^2}}
\end{align}

Under orthogonal rotation $R$ of the coordinate system:
\begin{align}
    \vec{C}'(h) &= R \vec{C}(h) \\
    \frac{d^2\vec{C}'}{dh^2} &= R \frac{d^2\vec{C}}{dh^2} \\
    \norm{\frac{d^2\vec{C}'}{dh^2}} &= \norm{R \frac{d^2\vec{C}}{dh^2}} = \norm{\frac{d^2\vec{C}}{dh^2}}
\end{align}

(since orthogonal matrices preserve norms).

Therefore, $v_0' = v_0$. $\square$
\end{proof}

\section{D.4 Fourier Analysis of Curvature}

\begin{theorem}
The Fourier components of $v_0(h)$ reveal structural periodicity at the helical repeat length.
\end{theorem}

\begin{proof}
DNA helical repeat: ${\sim}10.5$ bp/turn $\approx 35.7$\,\AA\ at 3.4\,\AA/bp

Computing FFT of $v_0(h)$:

\begin{align}
    \hat{v_0}(k) = \sum_h v_0(h) e^{-2\pi i k h}
\end{align}

Expected peak at:
\begin{align}
    k_{\text{peak}} = \frac{2\pi}{35.7 \text{ \AA}} = 0.176 \text{ \AA}^{-1}
\end{align}

Hypothesis: Structures with regular helical geometry show strong peaks; distorted structures show broadened peaks.

A full Fourier analysis of the kappa profiles across the 200-structure dataset is deferred to future work. The expected periodic signal at $k = 0.176$\,\AA$^{-1}$ (10-bp helical repeat) serves as an additional validation test.
\end{proof}

\section{D.5 Error Propagation}

\begin{theorem}
The uncertainty in $v_0$ propagates from coordinate uncertainty as:

\begin{align}
    \Delta v_0 \approx \frac{1}{(\Delta h)^2} \sqrt{3(\Delta x)^2 + 3(\Delta y)^2}
\end{align}

where $\Delta x, \Delta y$ are coordinate uncertainties from PDB B-factors.
\end{theorem}

\begin{proof}
Finite difference: $v_0 \approx \norm{[\Delta^2 x / (\Delta h)^2, \Delta^2 y / (\Delta h)^2]}$

Standard error propagation through finite differences:

\begin{align}
    \frac{\partial v_0}{\partial x_i} &\approx \frac{2\Delta x_i}{(\Delta h)^2} \\
    (\Delta v_0)^2 &= \sum_i (\frac{\partial v_0}{\partial x_i})^2 (\Delta x_i)^2
\end{align}

With $N=3$ terms contributing equally:

\begin{align}
    \Delta v_0 \approx \frac{1}{(\Delta h)^2} \sqrt{3(\Delta x)^2 + 3(\Delta y)^2}
\end{align}

For typical B-factors (20--40~\AA{}$^2$, uncertainty $\sim$0.5~\AA{}):

\begin{align}
    \Delta v_0 \approx \frac{\sqrt{3 \times 0.5^2 + 3 \times 0.5^2}}{1.0^2} \approx 0.01 \text{ L}^{-1}
\end{align}

$\square$
\end{proof}


